% Discussion section (5.1-5.7) of the journal paper; \input by main.tex.
%
% Built on the Results (results.tex) and on the conclusions and limitations
% of the final report (Final_correct_ results/IMD_Mapping_report 1.pdf,
% 2026-10-23). Interpretations beyond the report are worded as hypotheses.
% \litneeded{...} marks where published work still has to be found and cited;
% no reference has been added that is not already in refs.bib.

\section{Discussion}
\label{sec:discussion}

% ---------------------------------------------------------------------------
\subsection{Do AlphaEarth embeddings replace Sentinel-2?}
\label{sec:disc-replace}

The answer depends on how much clear Sentinel-2 imagery a city has. In
Milan, where 30 near cloud-free dates passed the screening, the Sentinel-2
stack and percentile composites outperformed AlphaEarth by 0.03--0.04 in
\Rsq{} and QWK, and both agreed with the photo-interpreted plots better than
CLMS itself (\cref{sec:res-milan}). In Hanoi and HCMC, where only three or
four dates passed a looser screening and the HCMC stack reduces to a single
image, AlphaEarth gave the best locally retrained map in both cities, ahead
of every Sentinel-2 composite (\cref{sec:res-hanoi,sec:res-hcmc}). Where
Sentinel-2 can be composed richly, it remains the stronger predictor; where
cloud cover limits it, the embeddings take the lead.

A likely reason is the content of the embeddings. Each AlphaEarth vector
summarises a full year of Sentinel-2, Sentinel-1, Landsat and other
observations \cite{brown2025alphaearth}, so it is not restricted to the few
clear optical dates that a local composite can use. The radar contribution in
particular is insensitive to cloud, which would explain why the embeddings
hold their accuracy in the monsoon climates of Hanoi and HCMC while the
Sentinel-2 composites lose theirs. This explanation is consistent with the
results but is not tested directly here: the study does not separate the
contribution of the individual sensors inside the embedding.

AlphaEarth was also the most stable predictor across settings. Its \Rsq{} for
local models ranged only from 0.590 to 0.617 across the three cities
(\cref{sec:res-cross}), it lost least when the training target was degraded
from CLMS to GHSL (+1.83\pp{} RMSE, against +3.31\pp{} for the stack;
\cref{tab:target-change}), and it needs no city-specific choice of dates,
cloud thresholds or compositing method. The same 64 bands are available for
every pixel and year, which makes a workflow built on them directly
portable. What the embeddings give up is the extra accuracy that a
well-composed, multi-date Sentinel-2 input achieves where the imagery allows
it.

The comparison between cross-validation and independent validation adds a
caution. In Milan AlphaEarth had the highest cross-validation error against
CLMS (\cref{tab:cv}) but ranked above the S2 median against the EarthLabel
plots (\cref{tab:milan-val}). Reproducing a reference product closely is not
the same as mapping imperviousness accurately, and predictors should be
compared against independent ground truth rather than against the product
used for training.

AlphaEarth is therefore not a general replacement for Sentinel-2 but a
strong default: it is the better choice where clear optical imagery is
scarce and the simpler one everywhere, while multi-date Sentinel-2
composites remain worth building where many clear dates exist.
\litneeded{published comparisons of AlphaEarth or other geospatial
foundation-model embeddings with Sentinel-2 composites, for imperviousness
or related continuous urban variables}

% ---------------------------------------------------------------------------
\subsection{Why local retraining matters}
\label{sec:disc-local}

Local retraining was the single largest improvement in the Vietnamese
cities. For AlphaEarth it raised \Rsq{} from 0.425 to 0.617 in Hanoi and from
0.224 to 0.605 in HCMC, and it reduced the bias from strong over-prediction
($-$16.2 and $-$21.7\pp) to moderate under-prediction (+6.0 and +11.1\pp;
\cref{tab:hanoi-val,tab:hcmc-val}). This gain was obtained with GHSL, a
target that in Milan performed clearly worse than CLMS
(\cref{sec:res-target}). A locally fitted model trained on an imperfect
global label thus outperformed a model trained on a better label in another
city.

The relationship between the predictors and imperviousness is evidently not
the same in Milan and in the Vietnamese cities. Several differences can
contribute: roofing materials and their spectral signatures, the density
and mix of buildings, roads and vegetation in the urban fabric, the
phenology of tropical rather than temperate vegetation, and the soil and
moisture conditions of pervious surfaces. A random forest learns the
mapping from predictors to imperviousness in the feature space it was
trained on and cannot adapt it to conditions it has not seen. Local training
points, even if labelled with GHSL, place the model in the right part of
that space.

Local retraining also requires little. The training points come from a
global product, GHSL, and the predictors from a global archive, so the
locally retrained workflow can be repeated for any city without local
reference data. Its remaining error is mostly the systematic
under-prediction inherited from the roof-only target (\cref{sec:disc-label}),
which is more tractable than the large and unpredictable offsets of
zero-shot transfer.
\litneeded{studies on the spatial transferability of random forest or other
machine-learning models for impervious surface or land-cover mapping across
cities or climate zones}

% ---------------------------------------------------------------------------
\subsection{Training-label quality versus predictor quality}
\label{sec:disc-label}

The Milan experiment isolates the training target and shows that it
matters more than the predictor. Replacing CLMS with GHSL lowered \Rsq{} by
0.057--0.101 for the same predictor, more than the largest difference
between any two predictors trained on the same target (0.054;
\cref{sec:res-target}). It also turned a small over-prediction into an
under-prediction of 7.7--8.3\pp. With GHSL as the target, the four
predictors performed almost identically (\Rsq{} 0.516--0.533): the target,
not the predictor, set the ceiling.

The problem with GHSL as a training target is a mismatch of definition
rather than random error. GHSL measures roofed built-up surface and excludes
roads and open pavements by design \cite{pesaresi2023ghsl}, whereas CLMS
\cite{clms2018imd} and the EarthLabel photo-interpretation both count every
sealed surface. A model trained on GHSL therefore learns that roads, car
parks and squares are pervious. This explains the spatial pattern of
\cref{fig:ghsl-clms}, with GHSL-trained maps below CLMS in the dense urban
core where such surfaces are common, and the consistent under-prediction of
GHSL itself by 17.6--19.0\pp{} in all three cities (\cref{tab:cross}).

GHSL-trained models nevertheless improved substantially on GHSL, raising
\Rsq{} from 0.234 to 0.516--0.533 in Milan and roughly halving the bias. A
plausible explanation is that a pixel-level regression cannot reproduce the
sharp building outlines of GHSL from 10\,m predictors. Roads and pavements
next to buildings share much of their spectral and contextual signature,
so the model assigns them part of the imperviousness of the buildings
around them. This partial recovery is a side effect, not a correction, and
it leaves a residual under-prediction of 6--11\pp{} in every GHSL-trained
map.

The practical consequence is that effort spent on the training target is
likely to return more than effort spent on predictors. In Milan, the
difference between the best and worst predictor was smaller than the cost
of the weaker label; in Vietnam, where only the weaker label exists, a
sealed-surface target comparable to CLMS would be expected to remove much
of the remaining bias. The residual bias of the GHSL-trained maps is also
systematic, and a small set of photo-interpreted plots, such as the
EarthLabel plots used here, could be used to estimate and correct it. That
correction was not tested in this study.
\litneeded{literature on label noise and label definition mismatch in
training data for remote-sensing regression; prior assessments of GHSL
against sealed-surface references}

% ---------------------------------------------------------------------------
\subsection{Why zero-shot transfer succeeds or fails}
\label{sec:disc-zeroshot}

Zero-shot transfer from Milan failed for AlphaEarth in both Vietnamese
cities and for the Sentinel-2 median in Hanoi. In every case the failure
was dominated by a level shift: the zero-shot maps lay 29--45\pp{} above
GHSL over the raster (\cref{tab:viet-diff}) and over-predicted the
photo-interpreted plots by 16--22\pp{}, except the S2 median in HCMC.

The shift has several possible sources. The first is a shift in the
predictors: the distribution of embedding values or reflectances in Hanoi
and HCMC differs from that in Milan, and a random forest, whose predictions
are averages of training values, maps unfamiliar inputs onto the Milan
cases they most resemble. The second is a shift in the relationship
itself, as discussed in \cref{sec:disc-local}: the same spectral or
embedding signature need not mean the same imperviousness in a temperate
European city and in a tropical Asian one. The third applies to Sentinel-2
only: the median model was trained on a composite of 30 Milan dates spread
over the year and applied to composites of three or four Vietnamese dates
from different seasons, so even identical surfaces would yield different
median values. AlphaEarth is free of this compositing mismatch, yet its
zero-shot maps were the worst in both cities. The shift therefore lies
mainly in the predictor--imperviousness relationship and not only in how
the inputs are composed.

The zero-shot maps nonetheless retained some information. The AlphaEarth
zero-shot map in Hanoi reached an overall accuracy of 0.827 at the 50\%
threshold, while placing only 3.3\% of the plots on their exact ordinal
level (\cref{tab:hanoi-val}). It separated mostly impervious from mostly
pervious plots, but at a systematically inflated level. This suggests that
a zero-shot model could be made useful by recalibrating its level with a
small number of local reference points, an option not evaluated here.

Transfer succeeded, in the sense of a high score, only for the Sentinel-2
median in HCMC. As the next section argues, that success does not reflect
a transferable model.
\litneeded{work on domain shift and domain adaptation in Earth-observation
regression, and on the transferability of foundation-model embeddings
across regions}

% ---------------------------------------------------------------------------
\subsection{The HCMC anomaly / apparent Sentinel-2 zero-shot advantage}
\label{sec:disc-hcmc}

In HCMC the Sentinel-2 median zero-shot map was the best of all maps on
every metric, with \Rsq{} 0.643 and QWK 0.782 (\cref{tab:hcmc-val}). Taken
alone, this result would suggest that Sentinel-2 transfers across
continents better than AlphaEarth and better than local retraining. The
Hanoi results rule this out. The same model, applied in the same way, was
the weakest of the five Hanoi maps on \Rsq{} (0.416).

The explanation is the level the model carries over from Milan
(\cref{fig:level}). At the validation plots it predicts a mean of 61.0\% in
Hanoi and 57.9\% in HCMC, while the photo-interpreted means are 44.4\% and
50.2\%. Its over-prediction is therefore 16.61\pp{} in Hanoi but only
7.64\pp{} in HCMC. The continuous and agreement metrics used here, \Rsq{}
computed against the 1:1 line, $\kappa$ and QWK, all penalise a level
offset, so a map whose offset happens to be small scores well. The locally
retrained HCMC maps, in turn, under-predicted by 10.8--11.4\pp, more than
in Hanoi, which narrows their advantage further. The HCMC ranking reflects a
coincidence of levels in one city, not mapping skill.

Two lessons follow for the evaluation of transfer. First, a transfer
experiment validated in a single target city can be misleading, because the
agreement between the source level and the target level is itself a matter
of chance; testing in more than one target city, as here, exposes it.
Second, the bias should be reported alongside agreement metrics, and
comparisons of the mean predicted level with the reference
(\cref{fig:level}) should accompany the error statistics, so that a good
score obtained through a matching level can be told apart from a good
map.

% ---------------------------------------------------------------------------
\subsection{Implications for operational IMD mapping where local reference
data are scarce}
\label{sec:disc-operational}

Hanoi and HCMC have no imperviousness product, and this situation is common
outside Europe and North America. The results point to a practical
workflow for such cities.

\begin{enumerate}
  \item \emph{Train locally rather than transfer.} A model trained in a
  well-mapped city should not be applied unchanged to a city in a different
  region and climate. Zero-shot transfer from Milan produced
  over-predictions of 16--22\pp{} and was unreliable even where it scored
  well (\cref{sec:disc-zeroshot,sec:disc-hcmc}).
  \item \emph{Use a global product as the training target, and expect its
  bias.} GHSL makes local retraining possible anywhere, and locally trained
  models reached \Rsq{} 0.55--0.62 in both Vietnamese cities. Because GHSL
  omits roads and pavements, the maps will under-predict sealed-surface
  imperviousness, by 6--11\pp{} in this study.
  \item \emph{Prefer AlphaEarth where clear imagery is scarce.} In cloudy
  tropical cities, AlphaEarth gave the best local maps without any
  selection of dates or compositing. Where many clear Sentinel-2 dates
  exist, a multi-date Sentinel-2 composite can add accuracy.
  \item \emph{Collect a small independent reference set.} The 450
  photo-interpreted EarthLabel plots per city were enough to rank the maps,
  expose the level matching in HCMC and quantify the bias of every map.
  Such a set costs far less than a full reference product and could also
  serve to correct the systematic bias of GHSL-trained maps
  (\cref{sec:disc-label}).
\end{enumerate}

The whole workflow, from sampling to tuning, prediction and export, runs in
Google Earth Engine \cite{gorelick2017gee}, so it can be repeated for other
cities without local computing infrastructure. The code and the predicted
maps are openly available \cite{brovelli2026zenodo}. For the purposes that
motivated this study, Local Climate Zone mapping and urban heat island
analysis, the locally retrained AlphaEarth maps are the most suitable
product available for Hanoi and HCMC, provided their under-prediction of
roads and pavements is taken into account.
\litneeded{the needs for imperviousness inputs in LCZ mapping and urban heat
island studies, and the availability of imperviousness products outside
Europe}

% ---------------------------------------------------------------------------
\subsection{Limitations}
\label{sec:limitations}

\paragraph{Label mismatch.} CLMS measures sealed surface and GHSL roofed
surface (\cref{sec:disc-label}). The Vietnamese models were trained on GHSL
only, so they learn a target that leaves out roads and pavements, and their
accuracy is bounded by that choice rather than by the predictors.

\paragraph{Temporal scope.} The study is restricted to 2018 and does not
address change over time. Whether the conclusions hold for other years,
and whether the embeddings or composites support consistent multi-year
monitoring, remains untested.

\paragraph{Geographic extent.} Three cities, one European and two
Vietnamese, are not a sufficient basis for general claims about
transferability. The zero-shot test runs in one direction only, from Milan
to Hanoi and HCMC, in a single year. The Sentinel-2 median zero-shot model
is also trained on a composite of 30 Milan dates but applied to composites
of only three or four dates in Vietnam.

\paragraph{Valid mask.} Training points and validation plots are drawn only
where every predictor has a value. Apart from water, identified with ESA
WorldCover \cite{zanaga2022worldcover}, the mask removes pixels that are
cloud or shadow on a stack date (up to 1.8\% of the study area). The maps
are not tested on those pixels, where the stack maps have no value.

\paragraph{Model and reference design.} A single model family, the random
forest \cite{breiman2001rf}, is used throughout; other regressors, or
models that use the spatial context of a pixel, could rank the predictors
differently. The random forest also averages training values in its
predictions, which pulls them away from the extremes of 0 and 100\%; this
is a likely reason why almost every trained map has a lower exact ordinal
agreement than the benchmark product of its city. The EarthLabel plot
value is quantised to ninths (steps of 11.1\pp), which sets a floor on the
continuous error that any map can reach. The ground truth is the finalised
dual-labeller interpretation, except for 22 replacement plots (6 in Milan,
16 in Hanoi) labelled jointly.

\paragraph{Statistical uncertainty.} With 450 plots per city, some of the
differences between maps are small: 0.008 in \Rsq{} between the stack and the
percentile composite in Milan, 0.013 in $\kappa$ between AlphaEarth and the
stack in Hanoi. No confidence intervals are reported. A bootstrap over the
plots would show which of these differences are robust, and the narrowest
ones should be read as ties.
